% The two-stage pipeline: factuality posteriors become coherence-MRF priors.
%
% Stage 2 is three boxes, not four: mining already emits a Level-1 coupling, so
% there is no separate compile step to draw. In the implementation
% `compile_sense` is called inside the miner (lcs/relation_miner.py) and a
% MinedRelation carries both its level2_sense and its level1_type; the scorer
% reads the coupling off the relation and never compiles anything. Drawing a
% separate "compile" box implied a stage that does not exist -- and overloaded
% $\mathcal{C}$, which is already the retrieved evidence set in Stage 1.
%
% The dashed arrow is labelled with the composition itself, $\pi_i \gets q_i$:
% the Stage-1 output is a POSTERIOR ($q_i$) and the Stage-2 unaries are the
% coherence MRF's PRIORS ($\pi_i$, exactly the $\phi_i=[1-\pi_i,\pi_i]$ of
% Def. def-mrf). The two symbols must stay distinct -- Alg. alg-lcs reports
% $\#\{q_i<\pi_i\}$, claims dragged below their own prior.
%
% Deliberately label-free of \ref{eq-priors}: submit.tex states the handoff
% inline and defines no such label, and this file is shared by both papers.
% ---------------------------------------------------------------------------
\begin{figure}[htbp]
\centering
\resizebox{\textwidth}{!}{%
\begin{tikzpicture}[
    font=\small,
    box/.style={draw,rounded corners=1.5pt,align=center,inner sep=3.4pt,
                minimum height=8.5mm},
    s1/.style={box,fill=blue!5},
    s2/.style={box,fill=orange!8},
    fin/.style={box,fill=green!10},
    ar/.style={-{Stealth[length=1.9mm]},thick},
    node distance=3.2mm]

  \node[s1] (resp) {response $y$};
  \node[s1,right=of resp] (atom) {atomize\\+ revise};
  \node[s1,right=of atom] (ret)  {retrieve\\contexts $\mathcal{C}$};
  \node[s1,right=of ret]  (nli)  {NLI\\$p^\ast$};
  \node[s1,right=of nli]  (fmrf) {factuality\\MRF};
  \node[fin,right=of fmrf] (post) {$q_i=\pr{a_i{=}1\mid\mathcal{C}}$};

  \draw[ar] (resp) -- (atom);
  \draw[ar] (atom) -- (ret);
  \draw[ar] (ret)  -- (nli);
  \draw[ar] (nli)  -- (fmrf);
  \draw[ar] (fmrf) -- (post);

  % Mining emits the coupling, so there is no separate compile box.
  \node[s2,below=10mm of atom] (mine)
       {mine relations\\sense $\to$ coupling\\(Prompt A / B)};
  \node[s2,right=of mine] (cmrf)
       {coherence MRF\\unaries $[1{-}\pi_i,\pi_i]$};
  \node[fin,right=of cmrf] (lcs) {$\LCS_{\mm}(y)$};

  \draw[ar] (atom.south) -- ++(0,-4mm) -| (mine.north);
  \draw[ar] (mine) -- (cmrf);
  \draw[ar] (cmrf) -- (lcs);
  \draw[ar,dashed] (post.south) -- ++(0,-5mm) -| (cmrf.north)
        node[pos=0.28,above,font=\scriptsize]{$\pi_i \gets q_i$};

  \node[left=1mm of mine,font=\scriptsize\itshape,align=right,text width=13mm]
       {Stage 2\\coherence};
  \node[left=1mm of resp,font=\scriptsize\itshape,align=right,text width=13mm,
        yshift=0mm] {Stage 1\\factuality};
\end{tikzpicture}}
% scripts/make_pipeline_png.sh defines \plnocaption to render this figure for
% standalone export, where surrounding prose supplies the explanation. Undefined
% in the paper, so the caption below is typeset normally there.
\ifdefined\plnocaption\else
\caption{The two-stage pipeline. Stage 1 (\S\ref{sec-factuality}) is FactReasoner
  \citep{factreasoner}: it grounds each atom against retrieved evidence and returns a
  posterior marginal $q_i$. Stage 2 (\S\ref{sec-coherence}) builds a second MRF over
  the atoms \emph{alone} and adds one pairwise factor per mined atom--atom relation;
  the mined relation already carries its Level-1 coupling, so no separate compile step
  intervenes. The dashed arrow is the only coupling between the stages, and carries
  the composition $\pi_i\!\gets\!q_i$ into the $\phi_i=[1-\pi_i,\pi_i]$ of
  Def.~\ref{def-mrf}. Atomization is shared, so the response is decomposed once and
  both stages score the same claim set; setting $\pi_i\!\equiv\!\tfrac12$ instead
  recovers a pure coherence score that consults no external evidence
  (Rem.~\ref{rem-noevidence}).}
\label{fig:pipeline}
\fi
\end{figure}
